\section*{Abstract}
Multimodal emotion recognition systems increasingly fuse their modalities with learned, attention-based modules rather than by averaging their predictions. This thesis asks when that added complexity pays off, using the EAV corpus: 42 participants recorded with 30-channel EEG, speech audio and facial video across five emotions. Per-modality encoders---the Audio Spectrogram Transformer, a facial-emotion Vision Transformer and EEGNet---feed four fusion methods: naive late fusion (mean softmax), a trimodal self-attention module, a concatenation MLP, and the attention module trained with softhard modality dropout, which also permits inference without EEG. Results are reported within-subject, as in all prior trimodal work on EAV, and under five-fold subject-wise cross-validation with every encoder retrained per fold.

Within-subject, the encoders reach 57.1\% (audio), 74.6\% (vision) and 43.9\% (EEG), and naive late fusion 77.5\%. Once model selection is kept off the test set, learned fusion is at best marginally better: 80.0\% for the concatenation MLP and 79.8\% for dropout-trained attention, while plain attention (76.8\%) does not beat averaging. On unseen participants the picture inverts. Naive late fusion reaches 62.2\%, every learned head is 3.4--4.5 points worse, and the loss is concentrated in vision, which falls by 33 points while audio transfers at parity.

The principal contribution is unsupervised subject calibration. Standardising each participant's fusion features with statistics from a short unlabelled recording lifts every learned head to 73.6--75.9\% on unseen participants, 11--14 points above averaging ($p_{\mathrm{Holm}} < 2\times10^{-7}$). Twenty unlabelled five-second clips suffice for a significant gain for every head, and the effect saturates at about fifty; a neutral-only enrolment session fails. Linear probes show that participant identity is decodable at 82--100\% from every modality before calibration and at chance after it. Under calibration the heads rank concatenation MLP $>$ dropout-trained attention $>$ attention, each step significant: attention-based fusion beats averaging decisively once identity is removed, and modality dropout is what makes it competitive.

The second contribution is an integrity audit of the thesis's own earlier results. Selecting fusion epochs on the test set had inflated them by 5.3--7.5 points, and the earlier state-of-the-art claim is withdrawn. Trial alignment across modalities cannot be verified from the released data. A suppression matrix previously read as population-level emotion suppression reproduces the EEG classifier's own errors ($r = 0.989$) and does not exceed a class-conditional permutation null.

\vspace{1cm}
\textbf{Keywords: Multimodal Emotion Recognition, Attention-Based Fusion, EEG, Subject-Independent Evaluation, Subject Calibration, Participant Identity, Modality Dropout, Evaluation Integrity, EAV Dataset, Affective Computing.}
\pagebreak
